% ===== 原文 PDF pp.79--82（印刷 pp.78--81）=====
% [边界] 本文件开头是 p.79 页首（印刷 p.78），承接 tr07a 的 p.78。末尾 p.82 的末句
%   （"…the Wess-Zumino contributions"）跨到 p.83，p.82 与 p.83 同属 tr-07，故该半句在
%   tr07c.tex 开头接写完（本文件末标 % [边界]）。
% [原文 p.78]

如第 2.3 节所述，一个 $\cN=2$ 理论拥有丰富的真空模空间；依据哪些 R 对称性子群被自发破缺，它可以被刻画为库仑分支、希格斯分支和/或混合分支。在 IIB 型紧化中，前两者分别与 CY 的复结构模和 Kähler 模相关。库仑分支上的低能物理尤其有趣，它由 Seiberg--Witten 几何 [20,21] 决定。研究一个 $\cN=2$ SCFT 最重要的任务之一，就是得到它的 SW 几何；在我们的情形里，这可以通过辨认 IHS 的 mini-versal 形变来实现\footnote{孤立奇点的形变理论在数学中是一个研究得很透彻的课题。请感兴趣的读者参阅经典文献 [272,273] 以了解精确的数学定义与定理。这里我们只需直观地指出：versal 形变是那些无法通过重定义形变参数而消去的奇点形变；而 mini-versal（极小 versal）形变则是参数个数最少的 versal 形变。物理上，mini-versal 形变给出相应 4d $\cN=2$ 理论库仑分支上的独立参数。}。

一个 IHS 的 mini-versal 形变可以如下描述：给定描述原点处 IHS 的拟齐次多项式 $F$，该 IHS 的 mini-versal 形变取下形式：
\begin{equation}
  \hat F = F(x_i) + \sum_{\alpha=1}^{\mu} g^{\alpha}(x_i)\,\lambda_{\alpha}\,.
  \tag{5.6}
\end{equation}
这里 $g^{\alpha}:=x_1^{\alpha_1}x_2^{\alpha_2}x_3^{\alpha_3}x_4^{\alpha_4}$（其中 $\alpha=(\alpha_1,\alpha_2,\alpha_3,\alpha_4)$）表示 Jacobi 代数（Milnor 环）$J(F)$ 的单项式基，而 $J(F)$ 定义为
\begin{equation}
  J(F) = \C[x_1,x_2,x_3,x_4]/(dF)\,,
  \tag{5.7}
\end{equation}
其中 $\C[x_i]$ 表示 $\C^4$ 的多项式环，$(dF)$ 表示由 $\partial_iF=0$ 生成的 Jacobi 理想（它消去平凡的形变）。Jacobi 代数的维数 $\mu$ 为\footnote{为推出这一点，例如可以构造代数 $J(F)$ 的 Poincaré 多项式并读出 $\mu$。细节可参见例如 [48] 及其所引文献。几何上，$\mu$ 也称为 Milnor 数，它计算形变后 Calabi--Yau $\hat W:=\{\hat F=0\}$ 的中间维数同调圈的个数。从 2d LG 的观点看，$\mu$ 与 Witten 指数 $\Tr(-1)^F$ 一致。}
\begin{equation}
  \mu = \prod_{i=1}^{4}\bigl(q_i^{-1}-1\bigr)\,.
  \tag{5.8}
\end{equation}

% [Tutor Note: (5.8) 的来历：对拟齐次 $F$，Jacobi 代数 $J(F)$ 的 Poincaré 多项式是
%   $\prod_{i=1}^4(1+t+\cdots+t^{k_i-1})$，其中 $q_i=1/k_i$（即 $F$ 是诸 $x_i^{k_i}$ 的准齐次组合）；
%   令 $t=1$ 即得 $\mu=\prod_i(1/q_i-1)$。自检：$A_k$ 奇点取
%   $q=(1/2,1/2,1/2,1/(k+1))$，得 $\mu=1\cdot1\cdot1\cdot k=k$，正是 $A_k$ 的 Milnor 数。]

% [Tutor Note: 例 5.1 的框体在原文中自 p.79 页末起（「Example 5.1:」+ 一句引导语），跨到 p.80 的
%   (5.9) 为止；本文件紧接 (5.8) 之后排版该框，顺序与原文一致。]

\begin{examplebox}{例 5.1：}
  我们给出一些由 IIB 型弦论探测 ADE 奇点所实现的 4d $\cN=2$ SCFT\footnote{注意，这类 4d SCFT 没有边缘形变与无关形变。}的例子，连同它们的 mini-versal 形变：
  % [原文 p.79]
  \begin{equation}
    \begin{aligned}
      F_{A_k}(x,y) &= x_1^2+x_2^2+x_3^2+x_4^{k+1}\,, &\quad
        \hat F_{A_k}(x,\lambda) &= F_{A_k}+\lambda_1+\lambda_2x_4+\cdots+\lambda_kx_4^{k-1}\\[2pt]
      F_{D_k}(x,y) &= x_1^2+x_2^2+x_3^{k-1}+x_3x_4^{2}\,, &\quad
        \hat F_{D_k}(x,\lambda) &= F_{D_k}+\lambda_1+\lambda_2x_3+\cdots+\lambda_kx_3^{k-1}\\[2pt]
      F_{E_6}(x,y) &= x_1^2+x_2^2+x_3^3+x_4^{4}\,, &\quad
        \hat F_{E_6}(x,\lambda) &= F_{E_6}+\lambda_1+\lambda_2x_3+\lambda_3x_4+\lambda_4x_4^2+\\[2pt]
      & & &\quad {}+\lambda_5x_3x_4+\lambda_6x_3x_4^2\\[2pt]
      F_{E_7}(x,y) &= x_1^2+x_2^2+x_3^3+x_3x_4^{3}\,, &\quad
        \hat F_{E_7}(x,\lambda) &= F_{E_7}+\lambda_1+\lambda_2x_3+\lambda_3x_4+\lambda_4x_4^2+\\[2pt]
      & & &\quad {}+\lambda_5x_4^3+\lambda_6x_3^2x_4+\lambda_7x_3x_4\\[2pt]
      F_{E_8}(x,y) &= x_1^2+x_2^2+x_3^3+x_4^{5}\,, &\quad
        \hat F_{E_8}(x,\lambda) &= F_{E_8}+\lambda_1+\lambda_2x_3+\lambda_3x_4+\lambda_4x_4^2+\\[2pt]
      & & &\quad {}+\lambda_5x_4^3+\lambda_6x_3x_4+\lambda_7x_3x_4^2+\lambda_8z_3z_4^3
    \end{aligned}
    \tag{5.9}
  \end{equation}
\end{examplebox}

% [Tutor Note: 例 5.1 的框体在原文中到 (5.9) 为止（p.80 的左边框在 (5.9) 下方收尾，
%   其后「The above $\hat F$ defines …」已回到正文）；此处从原文。]

% [Tutor Note: (5.9) 末行（$E_8$ 的形变）的最后一个单项式原文印作 $z_3z_4^3$。本列表中只有
%   $x_1,\dots,x_4$ 四个坐标，前面各项也都用 $x$，故此处 $z$ 应为 $x$：正确形式是
%   $\lambda_8x_3x_4^3$（500 dpi 裁图复核，原文确系印作 $z_3z_4^3$，非文字层误抽）。
%   正文照原文抄录，不改动。]

% [Tutor Note: (5.9) 的 $D_k$ 一行值得警惕。$F_{D_k}=x_1^2+x_2^2+x_3^{k-1}+x_3x_4^2$ 的 Jacobi
  %   理想由 $x_1,x_2,\,(k-1)x_3^{k-2}+x_4^2,\,2x_3x_4$ 生成，商环的单项式基是
  %   $\{1,x_3,\dots,x_3^{k-2},x_4\}$（$k$ 个单项式，与 $\mu(D_k)=k$ 相符），所以 mini-versal 形变应形如
  %   $\lambda_1+\lambda_2x_3+\cdots+\lambda_{k-1}x_3^{k-2}+\lambda_kx_4$。原文此处印作
  %   $\lambda_1+\cdots+\lambda_kx_3^{k-1}$（500 dpi 裁图复核确认如此印刷），与 Jacobi 基不符：
  %   $x_3^{k-1}=x_3\cdot x_3^{k-2}$ 在商环中为零。此处照原文抄录，不改动。]

  上述 $\hat F$ 定义了一个（广义的）Seiberg--Witten 几何：它是一个形变后的 CY 三重复形，纤维化在由 $\lambda_{\alpha}$ 参数化的特殊 Kähler 模空间之上，而这些 $\lambda_{\alpha}$ 正是要与 4d SCFT 中的 CB 数据相等同的量。这一构造与我们在前面几节所遇到的情形不同：在第 3.3 节中，SW 几何由一条 SW 曲线给出，即 $\hat F=\{(x,z,\lambda)\}$，纤维化在 CB 模空间之上，它一般来自 Class S 构造中的 Hitchin 系统。而在 IHS 上的 IIB 型几何构造中，SW 几何一般由一个三重复形（而非一条曲线）描述，纤维化在模空间之上 [26,31,268,274]。此外，SW 微分一形式被替换为 $\hat F$ 上的 Calabi--Yau 全纯典范三形式 $\Omega$\footnote{文献 [275] 中找到了它与所谓本原形式（primitive form）之间一个有趣的关联。}：
  \begin{equation}
    \Omega_3 = \frac{dx_1\wedge dx_2\wedge dx_3\wedge dx_4}{d\hat F(x_i,\lambda_{\alpha})}\,.
    \tag{5.10}
  \end{equation}

% [Tutor Note: 为什么 (5.10) 是这个形式？形变三重复形 $\hat F=0$ 的典范 $(3,0)$ 形式是 $\C^4$ 的
%   $\Omega_4=dx_1\wedge dx_2\wedge dx_3\wedge dx_4$ 沿法向的留数：把围道收在 $\hat F=0$ 上，
%   留数正是 $\Omega_4/d\hat F$。BPS 粒子的质量是 D3 膜缠绕 3-圈（特殊拉格朗日子流形）的
%   $\int\Omega_3$，其量纲为 1，恰与「质量」相符 —— 这也解释了下一段为何用 $\Omega_3$ 的荷来定 $a$。]

  所得 4d SCFT 的 CB 上的 BPS 态来自 D3 膜缠绕形变三重复形 $\hat F$ 上的超对称 3-圈（称为特殊拉格朗日子流形），它们的质量由 $\Omega_3$ 的积分决定\footnote{注意，许多 4d $\cN=2$ SCFT（尤其是由 II 型弦论实现的那些）与 2d $\cN=(2,2)$ SCFT 共享某些共同特征，包括 BPS 谱及其穿墙（wall-crossing）现象；这些细节我们在本综述中不展开，而请读者参阅 [72] 中关于所谓 4d--2d 对应的工作。}。只有在特殊情形下，(5.6) 中的 SW 几何才能约化为一条曲线纤维化，SW 微分 (5.10) 亦然。这样的例子出现在（非紧）CY 三重复形是 $\mathbb{P}^1$ 上的 ALE 纤维化之时 [26,164]。

  集合 $\{\lambda_{\alpha}\}$ 构成奇异的 Calabi--Yau 三重复形 $W$ 的复结构形变模，它们对应于所得 4d SCFT 库仑分支上保持 $\cN=2$ 超对称的那些参数，包括手征耦合、质量与库仑分支 VEV。要解读每个参数 $\lambda_{\alpha}$ 的物理含义，我们需要考察它们的标度维数 $\Delta(\lambda_{\alpha})$，或等价地考察它们的 $U(1)_r$ 荷。回忆第 2 节：4d SCFT 库仑分支上一个算子 $\cO$ 的标度维数 $\Delta(\cO)$（把它看作超手征主算子）正比于其 $U(1)_r$ 荷，而后者又正比于 $\C^*$ 作用下的荷 $q$，即 $\Delta(\cO)=aq(\cO)$，其中 $a$ 是比例常数。要确定 $a$，可以利用如下事实：典范三形式 $\Omega_3$ 在 $\C^*$ 下的荷为 $(\sum_{i=1}^{4}q_i-1)$，而它的标度维数为 $\Delta(\Omega_3)=1$——因为 $\Omega_3$ 在 3-圈上的积分给出一个 BPS 粒子的质量。于是有
  \begin{equation}
    a = \frac{1}{\sum_{i=1}^{4}q_i-1}\,.
    \tag{5.11}
  \end{equation}
  % [原文 p.80]
  现在我们要确定 (5.6) 中形变 $\lambda_{\alpha}g^{\alpha}$ 的参数 $\lambda_{\alpha}$ 的 $\C^*$ 荷 $q(\lambda_{\alpha})$。因为
  \begin{equation}
    q(\lambda_{\alpha}) = 1-q(g^{\alpha})\,,
    \tag{5.12}
  \end{equation}
  所以只需知道 $g^{\alpha}:=x_1^{\alpha_1}x_2^{\alpha_2}x_3^{\alpha_3}x_4^{\alpha_4}$ 的荷 $q(g^{\alpha})$，它是
  \begin{equation}
    q(g^{\alpha}) = \sum_{i=1}^{4}q_i\alpha_i\,.
    \tag{5.13}
  \end{equation}
  % [Tutor Note: (5.12) 之所以成立，是因为整个形变项 $\lambda_\alpha g^\alpha$ 与 $F$ 一样在 $\C^*$ 下荷为 1
  %   （见 (5.2)：$F(\xi^{q_i}x_i)=\xi F(x_i)$）。于是 $q(\lambda_\alpha)+q(g^\alpha)=1$。
  %   这一步也顺带说明：$\lambda_\alpha$ 的标度维数不是任意的，它被单项式的荷唯一决定。]
  于是我们得出结论：$\lambda_{\alpha}$ 的标度维数 $\Delta(\lambda_{\alpha})$ 由下式给出
  \begin{equation}
    \Delta(\lambda_{\alpha}) = \frac{1-\sum_{i=1}^{4}q_i\alpha_i}{\sum_{i=1}^{4}q_i-1}\,.
    \tag{5.14}
  \end{equation}

\begin{examplebox}{例 5.2：}
  考虑奇点 $F=x_1^{a_1}+x_2^{a_2}+x_3^{a_3}+x_4^{a_4}$，并带约束 $\sum_{i=1}^{4}\frac{1}{a_i}>1$。此时 Milnor 数为 $\mu=\sum_{i=1}^{4}(a_i-1)$。Jacobi 理想 $(dF)$ 由 $\partial_{x_i}F$ 生成，即 $(x_1^{a_1-1},x_2^{a_2-1},x_3^{a_3-1},x_4^{a_4-1})$。因此相应的 Jacobi 代数 $J$ 具有如下单项式基：
  \begin{equation}
    x_1^{\alpha_1}x_2^{\alpha_2}x_3^{\alpha_3}x_4^{\alpha_4}\,,\qquad 0\leqslant\alpha_i\leqslant a_i-2\,.
    \tag{5.15}
  \end{equation}
  % [原文 p.81]
  于是上述形变的系数的标度维数为
  \begin{equation}
    \Delta(\lambda_{\alpha}) = \frac{1-\sum_{i=1}^{4}\frac{\alpha_i}{a_i}}{\sum_{i=1}^{4}\frac{1}{a_i}-1}\,.
    \tag{5.16}
  \end{equation}
\end{examplebox}

形变参数 $\lambda_{\alpha}$ 的集合可以按其标度维数 $\Delta(\lambda_{\alpha})$ 分类。其中我们有：
\begin{enumerate}[label=\arabic*.,leftmargin=2.2em,itemsep=5pt,topsep=4pt]
  \item 标度维数 $\Delta(\lambda_{\alpha'})>1$ 的参数 $\lambda_{\alpha'}$ 是一个库仑分支 VEV 参数，即同标度维数的 CB 算子 $\cO_{\alpha'}$ 的 VEV $\langle\cO_{\alpha'}\rangle:=\lambda_{\alpha'}$。我们用 $r$ 表示库仑分支参数的个数，即该 4d $\cN=2$ SCFT 的秩。若存在弱耦合的规范理论描述，它将与规范群的秩一致。在这组 $\lambda_{\alpha'}$ 中，每个满足 $1<\Delta(\lambda_{\alpha'})<2$ 的参数 $\lambda_{\alpha'}$ 给出一个相关（relevant）的 CB 算子；类似地，若 $\Delta(\lambda_{\alpha'})=2$ 或 $\Delta(\lambda_{\alpha'})>2$，相应的 CB 算子分别是边缘的或无关的。
  \item 标度维数 $\Delta(\lambda_{\beta})<1$ 的参数 $\lambda_{\beta}$ 被认定为手征耦合常数，它伴随该 SCFT 的一个 $\cN=2$ 超对称 F-项形变 $\int \dd^4\theta\,\cO_{\alpha}\lambda_{\beta}$。特别地，每个标度维数 $\Delta(\lambda_{\beta})=0$ 的参数 $\lambda_{\beta}$ 对应于一个恰当边缘形变。
  \item 标度维数 $\Delta(\lambda_{\gamma})=1$ 的参数 $\lambda_{\gamma}$ 是质量参数，它可以被看作味对称群 $F$ 的极大 Cartan 环对应的背景矢量多重态中标量的 VEV。我们用 $f$ 表示质量参数的个数，即相应 4d $\cN=2$ SCFT 中味对称性 $F$ 的秩。
\end{enumerate}

由上面第二点可以断定：参数 $\lambda_{\alpha}$ 成对出现 $(\lambda_{\alpha'},\lambda_{\beta})$，满足
\begin{equation}
  \Delta(\lambda_{\alpha'})+\Delta(\lambda_{\beta}) = 2\,,
  \tag{5.17}
\end{equation}
质量参数是例外。每个这样的对 $(\lambda_{\alpha'},\lambda_{\beta})$ 构成一个保持 $\cN=2$ 的 F-项形变 $\int\dd^4\theta\,\lambda_{\beta}\cO_{\alpha'}$，其中 $\cO_{\alpha'}$ 是一个 CB 手征主算子，其 VEV 为 $\lambda_{\alpha'}$。于是由这种配对关系我们断定 $\mu=2r+f$，而维数 $\mu$ 可以等同于由味荷扩张后的电磁荷格的秩。

% [Tutor Note: 为什么是 $\Delta(\lambda_{\alpha'})+\Delta(\lambda_{\beta})=2$ 而不是 $3$？因为这里的形变是写成
%   $\int\dd^4\theta\,\lambda_\beta\cO_{\alpha'}$ 的 F-项：4d 中 $\dd^4\theta$ 的标度维数是 2，
%   于是该项恰好边缘（marginal）的条件正是两者维数之和为 2。用 CB 的语言说，每个 $\Delta>1$ 的 VEV
%   参数都配一个 $\Delta<1$ 的手征耦合，后者把前者的 VEV 变成超势里的一个边缘项；质量参数
%   （$\Delta=1$）自配自，故不在配对计数之内，这正是 $\mu=2r+f$ 的来历。]

给定一个 4d $\cN=2$ SCFT 的完整参数谱 $\lambda_{\alpha}$，人们可以进而计算第 2.6 节引入的共形中心荷 $(a,c)$——它们度量该理论的自由度。考虑式 (2.50)，Wess--Zumino 贡献

% [边界] 本句（"Considering Eq. (2.50), the Wess-Zumino contributions R(A) and R(B) are
%   determined ..."）自 p.82 页末起、到 p.83 页首结束；p.82 与 p.83 同属 tr-07，
%   故此句在 tr07c.tex 开头接写完（不重复），并在那里标 % [原文 p.82]。
